% SGA 3，第 VI_B 篇，第一轮英文重建稿的中文翻译。
% 范围：段落10.14、定义10.15、命题10.16及其完整证明，含编者注102--103。
% 依据：Polo--Gille Exp6B-13oct24.pdf，本组件局部页74--75
% （合卷阅读版页415--416；再版页眉404--405；原版页边印刷页394--395）。
% 第11节始于局部页76，不在本组件内。
% OCR及Reinhold英文Markdown仅用作对照见证。
% 本组件不含图。

\medskip
\phantomsection\label{par:sga3-VIB-10.14}%
\noindent\textbf{10.14.---}
关于有限呈示的 \(S\)-群范畴的论断
\hyperref[lem:sga3-VIB-10.2]{10.2} 和
\hyperref[lem:sga3-VIB-10.3]{10.3}，可推广到如下的配对所成的范畴：
一个有限呈示的 \(S\)-群，以及一个以 \(G\) 为算子群的有限呈示
\(S\)-概形。更确切地，在 \hyperref[par:sga3-VIB-10.0]{10.0}
开头回顾的情形下：

\begin{enumerate}
\item[\textup{(i)}]\phantomsection\label{concl:sga3-VIB-10.14-i}
  设 \(j\in I\)，\(G_j\) 和 \(G'_j\) 是两个有限呈示的 \(S_j\)-群；
  设 \(H_j\)\textup{（分别地，\(H'_j\)）}是以 \(G_j\)
  \textup{（分别地，\(G'_j\)）}为算子群的有限呈示 \(S_j\)-概形。
  对 \(i\in I\)、\(i\geq j\)，令
  \[
    G_i=G_j\times_{S_j}S_i,
    \qquad
    G=G_j\times_{S_j}S,
  \]
  并类似地定义 \(G'_i\)、\(G'\)、\(H_i\)、\(H\)、\(H'_i\) 和 \(H'\)。
  以
  \[
    \operatorname{Dihom}_{S\text{-}\mathrm{gr.}}
      \bigl((G,H),(G',H')\bigr)
  \]
  表示由配对 \((G,H)\) 到配对 \((G',H')\) 的 \(S\)-群与带算子群的
  \(S\)-概形的双态射集合。则典范映射
  \phantomsection\label{form:sga3-VIB-10.14-Dihom-limit}%
  \[
    \varinjlim_{i\geq j}
      \operatorname{Dihom}_{S_i\text{-}\mathrm{gr.}}
        \bigl((G_i,H_i),(G'_i,H'_i)\bigr)
      \longrightarrow
      \operatorname{Dihom}_{S\text{-}\mathrm{gr.}}
        \bigl((G,H),(G',H')\bigr)
  \]
  是双射。

\item[\textup{(ii)}]\phantomsection\label{concl:sga3-VIB-10.14-ii}
  设 \(G\) 是有限呈示的 \(S\)-群，\(H\) 是以 \(G\) 为算子群的
  有限呈示 \(S\)-概形。则存在指标 \(j\in I\)、有限呈示的
  \(S_j\)-群 \(G_j\)、以 \(G_j\) 为算子群的有限呈示 \(S_j\)-概形
  \(H_j\)，以及由
  \[
    \bigl(G_j\times_{S_j}S,\ H_j\times_{S_j}S\bigr)
  \]
  到 \((G,H)\) 的 \(S\)-群与带算子群的 \(S\)-概形的双同构。%
  \begingroup
  \renewcommand{\thefootnote}{\fnsymbol{footnote}}%
  \footnote[1]{\label{note:sga3-VIB-10.14-GH}
  译者注：法文再版此处印作 \((G,X)\)。由于本断言以配对 \((G,H)\)
  开始，且 \SGARefLink{par:sga3-VIB-10.14}{10.14} 中并未定义 \(X\)，
  现恢复为 \(H\)。}%
  \endgroup
\end{enumerate}

\medskip
\phantomsection\label{def:sga3-VIB-10.15}%
\noindent\textbf{定义 10.15.---}%
\footnote{\label{note:sga3-VIB-102}编者注：通过区分“形式齐性”对象和
“齐性”对象的概念，更正了原文；参见IV，\S\S
\SGARefLink{sga3:iv:subsection:5:1}{5.1}和\SGARefLink{sga3:iv:subsection:6:7}{6.7}。}
\emph{设 \(\mathcal T\) 是 \((\mathbf{Sch}/S)\) 上比典范拓扑粗的拓扑。
给定 \(S\)-群概形 \(G\) 及以 \(G\) 为算子群的 \(S\)-概形 \(X\)，
若对每个 \(S'\to S\)、\(g\in G(S')\) 和 \(x\in X(S')\) 定义的态射}
\phantomsection\label{form:sga3-VIB-10.15-Phi}%
\[
  \Phi:G\times_S X\longrightarrow X\times_S X,
  \qquad
  (g,x)\longmapsto(gx,x),
\]
\emph{是拓扑 \(\mathcal T\) 的层范畴中的满态射，则称 \(X\) 为
\emph{\(G\) 下的形式齐性空间}\textup{（相对于拓扑 \(\mathcal T\)）}；
等价地，\(\Phi\) 是拓扑 \(\mathcal T\) 的覆盖
\textup{（参见 \hyperref[sga3:iv:proposition:4:4:3]{IV，4.4.3}）}。}

\emph{若 \(X\) 形式齐性，且态射 \(X\to S\) 也是拓扑
\(\mathcal T\) 的覆盖，则称它为\emph{齐性空间}。}

\emph{特别地，若 \(\Phi\) 是同构，则称 \(X\) 为
\emph{\(G\) 下的形式主齐性空间}；若 \(\Phi\) 是同构，且态射
\(X\to S\) 是拓扑 \(\mathcal T\) 的覆盖，则称 \(X\) 为
\emph{主齐性丛}\textup{（或 \(G\)-挠子）}
\textup{（参见 \hyperref[sga3:iv:definition:5:1:5]{IV，5.1.5}和
\hyperref[sga3:iv:proposition:5:1:6:item:ii:01]{5.1.6(ii)}）}。}

\medskip
\phantomsection\label{prop:sga3-VIB-10.16}%
\noindent\textbf{命题 10.16.---}
\emph{沿用 \hyperref[par:sga3-VIB-10.0]{10.0} 开头所考虑的情形。
设 \(j\in I\)，\(G_j\) 是 \(S_j\)-群，\(X_j\) 是以
\(G_j\) 为算子群的 \(S_j\)-概形。假设 \(G_j\) 和 \(X_j\)
在 \(S_j\) 上有限呈示。}

\emph{使}
\[
  X=X_j\times_{S_j}S
\]
\emph{成为 \textup{(fppf)} 拓扑下}
\[
  G=G_j\times_{S_j}S
\]
\emph{的齐性空间\textup{（分别地，主齐性丛）}的充要条件是存在指标
\(i\geq j\)，使}
\[
  X_i=X_j\times_{S_j}S_i
\]
\emph{是}
\[
  G_i=G_j\times_{S_j}S_i.
\]
\emph{的齐性空间\textup{（分别地，主齐性丛）}。}

结合 \hyperref[par:sga3-VIB-10.14]{10.14} 及 EGA~IV\(_3\)，
8.8.2、8.8.3和8.10.5，所述论断由
\hyperref[par:sga3-VIB-10.1]{10.1} 中回顾的 \textup{(fppf)} 拓扑的
覆盖态射的性质推出。%
\footnote{\label{note:sga3-VIB-103}编者注：关于用 \v{C}ech 上同调表达的
挠子的极限过渡性质，另参见SGA~4，VII，5.7及其推论。
\textup{[Ma07]}中对此有详细论述。}

% 中文翻译检查点：本组件全部正文、证明、脚注及源注已译；源155行。
% 机械核对：行内数学81处、陈列数学9处及全部引用键守恒（数学中的自然语言已译）。
% 状态：完成源文覆盖和逐段自校；尚未经独立合卷编译与页面审查。
